\documentclass[conference]{IEEEtran}
\IEEEoverridecommandlockouts
\usepackage{cite}
\usepackage{amsmath,amssymb,amsfonts}
\usepackage{booktabs}
\usepackage{array}
\usepackage{url}
\usepackage{graphicx}
\usepackage{textcomp}
\usepackage{xcolor}
\def\BibTeX{{\rm B\kern-.05em{\sc i\kern-.025em b}\kern-.08em
    T\kern-.1667em\lower.7ex\hbox{E}\kern-.125emX}}

\begin{document}

\title{Digital Memory Beyond the Upload Metaphor: An Engram Fidelity Ladder and Causal Transfer Atlas}

\author{\IEEEauthorblockN{Four-Agent Research Pipeline}
\IEEEauthorblockA{Survey--Idea--ExperimentDesign--Author Workflow\\
Research proposal; no experiments executed}}

\maketitle

\begin{abstract}
The prospect of storing, manipulating, and transplanting human memory in a digital system is often described as a future upload problem. That metaphor hides the scientific question: which properties of a memory must a representation preserve before it can be called a memory rather than a decoder, a behavioral policy, or a neural correlate? This paper proposes the Digital Engram Fidelity Ladder and Causal Transfer Atlas (DEFL-CTA), a design-first framework that evaluates memory representations at five explicitly separated levels. L0 measures predictive access, L1 behavioral association, L2 neural reinstatement, L3 content-preserving reconstruction, and L4 transfer or recollection under independent neural, behavioral, and content tests. The framework is grounded in evidence from human memory neuroscience, rodent engram manipulation, and human neural prostheses, while refusing to equate any one of these results with complete digital transplantation. We formalize a distributed, context-dependent state-space model; define content, relational, omission, false-detail, and collateral-memory endpoints; and specify held-out cue, cross-session, and cross-subject tests. The resulting proposal makes an ambitious claim: useful digital memory technology can be developed incrementally, but the decisive frontier is not higher decoding accuracy alone. It is selective, content-preserving, causally validated transfer that remains distinguishable from nuisance state, rehearsal, generic arousal, and confabulation. The protocol is design-only and reports no observed upload, editing, stimulation, or transplantation.
\end{abstract}

\begin{IEEEkeywords}
human memory, engrams, neural prostheses, digital representation, causal transfer, false memory, neurotechnology
\end{IEEEkeywords}

\section{Introduction}
Science fiction asks whether a person could upload a lifetime of memories to a cloud, edit a painful event, or download an episode into another body. The question is scientifically productive precisely because it forces several ordinarily conflated achievements apart. A system may predict which item a participant will recognize without containing the item's experiential content. A stimulation pattern may evoke a learned behavioral response without transferring an episode. A generative model may produce a fluent account that is semantically plausible but includes unsupported details. Conversely, a useful memory technology need not reproduce an entire person to preserve a well-defined event, relation, or skill.

Human memory is not a single file stored at one address. Cognitive systems distinguish episodic, semantic, procedural, and affective forms, and their expression depends on encoding, consolidation, retrieval, reconsolidation, context, and current state \cite{Squire2011,Kandel2014}. Psychological evidence also shows that remembering is reconstructive: omission, source confusion, intrusion, and confidence errors can coexist with a correct gist \cite{Schacter1999}. A digital representation therefore requires a declared target and a declared error model. Without those declarations, ``memory upload'' is not a testable scientific claim.

At the biological level, engram research has demonstrated that manipulating tagged neural ensembles can alter memory-related behavior in rodents \cite{Josselyn2020,Liu2012,Ramirez2013,Chen2019}. At the human interface level, neural prosthetic and closed-loop stimulation studies have reported assistance in constrained memory tasks \cite{Hampson2018,Ezzyat2018}. These are important causal and translational advances. They do not by themselves establish that a complete human episodic memory has been digitized, edited at will, or transplanted into an independent recipient. The gap is not a semantic technicality: it determines the measurements, controls, safety obligations, and conclusions that a valid study may support.

This paper presents DEFL-CTA, a research program selected from a structured Survey and Idea process. Its central thesis is that the field should replace the binary question ``Can memory be uploaded?'' with a promotion ladder. Each promotion requires new evidence, and no higher level can be inferred from a lower-level success. The ladder is paired with an atlas of causal interventions and failure modes so that a representation is evaluated for what it preserves, what it changes, and what it leaves unknown.

The contributions are fourfold. First, we define five operational fidelity levels and a phenotype ontology that separates content, temporal order, source context, relations, neural reinstatement, behavior, and subjective recollection. Second, we propose a distributed dynamical model whose latent variables can be tested against held-out cues while controlling for state, movement, arousal, and subject identity. Third, we specify a staged storage, manipulation, and transfer protocol with independent content scoring and collateral-memory endpoints. Fourth, we provide conditional decision rules that allow strong progress without overstating what a result means. The complete artifact remains a proposal: no human or animal data were executed in this run.

\section{Background, Survey, and Research Gap}
\subsection{Memory as a distributed and reconstructive process}
The computational analogy of input, encoding, storage, and retrieval is useful, but a literal file analogy is not. Memory can be expressed through several interacting systems and can change when it is retrieved. The neural substrate is distributed over hippocampal, cortical, and related networks, with different temporal roles across learning and recall \cite{Squire2011,Kandel2014}. The same event can therefore be represented by multiple features: an item, its temporal order, the place and source in which it was encountered, its relation to other items, and its affective or action value.

The relevant object for a digital study is consequently a task-defined memory phenotype. Let $m$ denote an event or learned relation, and let $\Phi(m)$ denote a vector of prespecified features. A phenotype is not a claim that all features exhaust the original experience. It is a contract that states what the experiment can measure. For example, a paired-associate task may define item identity and pair relation, while an episodic task may additionally define temporal order and source context. A fear-conditioning task may define cue-specific defensive behavior, but that endpoint must not be renamed subjective recollection.

The reconstructive character of memory changes how a digital artifact should be scored. A fluent narrative is not automatically faithful; it may contain semantic completion or model-generated detail. A confident report is not automatically a ground truth; confidence and source monitoring can dissociate. The survey therefore requires independent event records, held-out probes, confidence calibration, and explicit omission and false-detail measures. This makes errors visible instead of allowing a single average accuracy to conceal them.

\subsection{Engrams and causal manipulation}
An engram is best treated as an experimentally defined memory-related neural ensemble or state, not as a universal synonym for a complete memory. Optogenetic work in rodents showed that stimulation of tagged hippocampal ensembles can activate fear-memory-related behavior \cite{Liu2012}. Related work used hippocampal manipulation to create a false-memory-like behavioral expression, demonstrating that ensemble-level interventions can influence what an animal does in a constrained context \cite{Ramirez2013}. Chen \textit{et al.} reported artificial enhancement and suppression of hippocampus-mediated memories, providing a further example of causal influence on memory-related expression \cite{Chen2019}.

These findings establish a valuable principle: neural representations can be causally relevant to memory expression. They also establish the boundary of the inference. A rodent's freezing response does not independently validate the content, temporal structure, or subjective recollection of a human episode. Ensemble tagging can be sparse, task-specific, state-dependent, and affected by the chosen behavioral assay. DEFL-CTA uses engram studies as a mechanistic bridge for causal participation, not as proof that the brain contains a portable binary object.

\subsection{Human neural assistance}
Human neural prosthetic research has moved the question from philosophical speculation toward measurable intervention. Hippocampal recording and stimulation have been used to develop models intended to facilitate encoding and recall in constrained tasks \cite{Hampson2018}. Closed-loop temporal-cortex stimulation has also been associated with improved memory performance under task-specific conditions \cite{Ezzyat2018}. These studies motivate a realistic engineering sequence: first identify task-relevant states, then design a subject-specific policy, then test whether intervention effects are selective and reproducible.

Their results are not equivalent to general digital storage. A prosthesis may provide a cue, compensate for a disrupted computation, or improve a participant's strategy. A decoder may infer an upcoming response while the underlying content remains inaccessible. A stimulation policy may facilitate an existing memory rather than deliver a memory to a new subject. The proposal therefore reports assistance, decoding, and reinstatement separately, and it requires cross-context content tests before using the language of transfer.

\subsection{Survey gap ledger}
The Survey Agent organized evidence into subhypotheses and accepted six high-confidence gaps. G1 is the lack of a unified operational definition of memory fidelity. G2 separates causal engram participation from portable content representation. G3 separates human prosthetic assistance from general digital storage. G4 concerns the distributed and dynamic character of representations across memory phases. G5 concerns the measurement of omission and confabulation. G6 concerns collateral effects, agency, autonomy, and consent. Cross-subject identity-preserving transfer and a causal equivalence criterion were retained as candidate gaps rather than established facts.

The evidence policy is deliberately asymmetric. A source can support a bounded claim about the paradigm it studied, but a candidate paper cannot be promoted to direct evidence merely because it was found through citation expansion. In the present registry, E1--E10 are the allowed sources. The registry is summarized in Table~\ref{tab:evidence} and retained in the machine-readable Survey manifest. The verified evidence supports an incremental research program; it does not support a claim that a human memory upload has already occurred.

\begin{table}[t]
\caption{Evidence registry and admissible interpretation}
\label{tab:evidence}
\centering
\footnotesize
\setlength{\tabcolsep}{2pt}
\begin{tabular}{p{0.08\columnwidth}p{0.38\columnwidth}p{0.43\columnwidth}}
\toprule
ID & Evidence family & Interpretation permitted by this proposal \\
\midrule
E1,E8 & Human memory systems and molecular/systems biology & Memory is heterogeneous, distributed, and dynamic. \\
E2 & Engram review & Engram evidence must be tied to a defined task and endpoint. \\
E3--E5 & Rodent ensemble manipulation & Causal influence on constrained memory-related behavior. \\
E6,E7 & Human prosthesis and closed-loop stimulation & Task-specific assistance is feasible; general storage is open. \\
E9 & Reconstructive memory and errors & Omission, intrusion, source error, and confidence are required endpoints. \\
E10 & Neurotechnology ethics & Governance, consent, agency, and autonomy are design variables. \\
\bottomrule
\end{tabular}
\end{table}

\section{Research Questions and Planned Contributions}
The primary research question is: can a content-specific memory state be represented digitally with enough fidelity that it can be manipulated, transferred, and independently recognized after reconstruction, while preserving the distinction between neural correlation, behavioral expression, and subjective recollection?

This question is decomposed into five operational questions:
\begin{enumerate}
\item Which features of a task-defined memory generalize across cues, sessions, and contexts?
\item Does a distributed latent state predict content-specific reinstatement beyond stimulus, motor, arousal, and subject-state variables?
\item Can a stored representation reconstruct content and relations while bounding omissions and false details?
\item Can a targeted digital edit change one memory dimension without changing neighboring memories, agency, or general cognition?
\item What evidence is sufficient to promote a system from decoder or association to content-preserving reconstruction and transfer?
\end{enumerate}

The planned contributions are not a single upload device. They are a reproducible measurement infrastructure: a memory phenotype schema; the L0--L4 fidelity ladder; a distributed state-space emulator; a versioned digital object that records uncertainty and unknown details; and a causal transfer atlas that maps interventions to target and collateral outcomes. This infrastructure can support future clinical and neuroprosthetic work while making a strong scientific claim possible: digital memory research becomes meaningful when content-level identity survives held-out tests, not when a model merely predicts behavior.

\section{Idea Source Checkpoints and Direction Selection Audit}
The Idea Agent began with the accepted Survey gaps and generated multiple independent routes. R1 proposed a memory phenotype ontology and fidelity ladder. R2 proposed a distributed engram state-space emulator. R3 proposed cross-subject content alignment. R4 proposed a closed-loop memory prosthesis. R5 combined the strongest components into DEFL-CTA.

The selection was based on coverage, falsifiability, independence, and feasibility. R1 directly resolves G1 and G5 but does not by itself model mechanism. R2 addresses G2 and G4 but could remain a decoder if content probes are absent. R3 directly challenges G7 but carries high risk of subject-specific semantic and experiential confounds. R4 addresses G3 and G6 but can be mistaken for assistance rather than transfer. R5 was selected because it makes each route a gate in one auditable sequence and preserves the failures rather than averaging them away.

The search rejected a pure ``memory upload'' route because it assumes the conclusion: it treats memory as a stable file before specifying its content. It also rejected decoder-only transplantation because a decoder or stimulation policy may be portable while the underlying event is not. A pure rodent fear-engram route was retained as a mechanistic component but rejected as the full direction because it lacks human episodic content and an independent digital transfer criterion.

The evolution checkpoints were: (N0) separate content, behavior, neural state, and recollection; (N1) add held-out content and relational reconstruction; (N2) represent memory as a time-dependent distributed state; and (N3) add collateral and selective-causal promotion gates. These edits convert an evocative question into a sequence of increasingly demanding tests without precluding a technically ambitious outcome.

\section{Problem Definition, Assumptions, and Hypotheses}
\subsection{Definitions}
For this proposal, a digital memory representation is a versioned object containing a task schema, feature and relation definitions, learned parameters, uncertainty, cue-to-state mappings, provenance, and explicitly unobserved details. It is not a generated narrative alone. Storage means that the object supports prespecified held-out queries after serialization and reloading. Manipulation means that a defined target dimension is changed under an intervention while non-target dimensions are measured. Transfer means that an independent recipient system uses the representation under novel cues and passes content, neural, behavioral, and collateral tests.

The fidelity levels are:
\begin{itemize}
\item L0, predictive access: predict recall, recognition, or item identity on held-out trials.
\item L1, behavioral association: produce a target-related response in a constrained task.
\item L2, neural reinstatement: reproduce a target-specific neural state under held-out cues.
\item L3, content-preserving reconstruction: recover content, temporal order, and relations with bounded error.
\item L4, transfer or recollection: an independent recipient expresses target content under novel probes and passes neural, behavioral, and independent content checks.
\end{itemize}

\subsection{Assumptions}
The protocol assumes that a memory target can be made measurable without claiming to exhaust experience. It assumes that trial-level stimulus, context, response, confidence, and timing metadata can be obtained under an approved data-use protocol. It assumes that model capacity and preprocessing can be locked before evaluation, and that subject and session splits can be made before feature learning. It does not assume that neural similarity equals content identity, that cross-subject alignment is unique, or that subjective recollection can be directly observed from a neural trace.

\subsection{Hypotheses}
H1 states that distributed context-dependent latent states outperform stimulus-only and nuisance-only models on held-out content and neural reinstatement, after controlling for movement, arousal, task, and history. H2 states that L3 reconstruction requires joint preservation of item identity, temporal order, and relational structure; a system with high item accuracy but high false-detail or relation error will fail L3. H3 states that targeted edits will show a selective causal effect on the target dimension while neighboring memories and general task performance remain within preregistered bounds. H4 states that cross-subject transfer will be weaker than within-subject transfer unless alignment preserves content and relations rather than only semantic or motor policy. H5 states that confidence and narrative fluency will not be treated as sufficient evidence of fidelity and may increase while false details also increase.

\section{Study Design and Methods}
\subsection{D0: Fidelity ontology and preregistration}
Before any analysis, the team preregisters the memory class, task, cue set, encoding and retrieval windows, neural modality, inclusion rule, context variables, data splits, model capacity, primary endpoints, and manipulation targets. Primary endpoints are assigned by level rather than post hoc. For L0, the endpoint is held-out predictive performance. For L1, it is target-specific behavioral association. For L2, it is neural pattern similarity and temporal generalization. For L3, it is a vector of content and relation scores with error bounds. For L4, all of these are required in a recipient setting.

The preregistration also defines promotion and failure rules. Null results remain in the atlas. If a model violates the data contract or a batch fails its output schema, it is discarded as an invalid batch rather than silently converted into evidence. Missing labels, poor synchronization, or insufficient repeated trials become \texttt{needs\_human\_input}. This prevents a pipeline failure from being mistaken for a biological failure and prevents an unmeasured feature from being counted as preserved.

\subsection{D1: Data contract and structured memory task}
The preferred source is an ethically approved, de-identified human dataset with trial-level stimuli, item and relation labels, response, confidence, response time, and neural recordings such as intracranial electroencephalography, scalp electroencephalography, magnetoencephalography, functional magnetic resonance imaging, or a multimodal combination. If an approved animal engram dataset is used, cell-tagging, context, stimulation condition, and the memory-related assay must be retained. Data provenance includes subject, session, task, modality, preprocessing, electrode or region metadata, and consent or protocol identifiers.

The task should contain item identity, temporal order, source context, and relational attributes. Repeated encoding and delayed retrieval separate learning, consolidation, retrieval, and reconsolidation. Novel probes combine features not presented together as a single cue. Negative controls include non-target memories, semantically related foils, state-matched non-memory trials, and motor- or arousal-matched trials. The analysis is split by subject and session before feature learning. Within each training fold, preprocessing and hyperparameters are estimated only from training data.

\subsection{D2: Distributed engram state-space model}
Let $z_t$ be a latent memory-related state, $x_t$ a sensory or task input, $c_t$ the context, and $q_t$ measured nuisance and history covariates. A candidate dynamical model is
\begin{equation}
z_t=A_{c_t}z_{t-1}+B_{c_t}x_t+\epsilon_t,
\label{eq:state}
\end{equation}
where $\epsilon_t$ represents process uncertainty. For spiking data, the observation model is
\begin{equation}
\lambda_i(t)=\exp\left(b_i+w_i^\top z_t+v_i^\top q_t\right),
\label{eq:observation}
\end{equation}
where $\lambda_i(t)$ is the conditional intensity for channel $i$. For imaging or field-potential modalities, the point-process likelihood is replaced by a modality-appropriate observation model while preserving the latent-state and context separation.

Five nested model families are compared: stimulus-only, context-only, memory-state, memory-state-plus-history, and distributed memory-state. Evaluation uses held-out predictive log likelihood, calibration, neural pattern similarity, temporal generalization, and content probes. The model is a measurement device, not a claim that the brain literally implements a linear system. A candidate memory state must predict target-specific held-out responses beyond state, motor, arousal, and session controls. Subject and session leakage is assessed by shuffled labels and a subject-shuffle baseline.

\subsection{D3: Digital storage and reconstruction}
The stored object contains the task schema, learned state-space parameters, uncertainty, cue-to-state map, feature relations, provenance, and an explicit list of unknown details. Four reconstruction systems are compared: a behavioral classifier, a neural-state decoder, a generative latent-state model, and a hybrid model. All systems receive the same preregistered train/test partitions. A generated narrative may be used as a display layer, but its details are scored against independently held-out event records; unsupported plausibility is an error.

For a target memory $m$, define the content score as
\begin{equation}
F_{\mathrm{content}}(m)=1-\frac{E_{\mathrm{omission}}(m)+E_{\mathrm{false}}(m)+E_{\mathrm{relation}}(m)}{Z_m},
\label{eq:content}
\end{equation}
where omission, false-detail, and relational errors are independently scored and $Z_m$ is a preregistered normalization. Report the components, not only their aggregate. Additional endpoints include item accuracy, temporal-order accuracy, source-context accuracy, relation consistency, neural reinstatement similarity, response time, confidence calibration, and performance on novel cue combinations. A high behavioral score cannot hide a high false-detail rate.

\subsection{D4: Digital manipulation and collateral effects}
The first intervention is in silico. One target feature, temporal relation, or contextual attribute is changed while other dimensions are held fixed. The system predicts target recall, neighboring-memory responses, affect, decision confidence, and neural reinstatement. Matched edits to a non-target memory, sham edits, and state-only edits establish specificity. The primary contrast is target change minus nuisance-matched and non-target change.

If a future approved neural intervention is feasible, a model-defined dimension may be translated into a subject-specific modulation policy only after safety, feasibility, and autonomy review. The design requires target, sham, non-target, and nuisance-matched perturbation conditions. Outcomes include target memory, neighboring memories, source-monitoring errors, false details, confidence, agency, affect, and task strategy. A global recall increase, generalized defensive response, or altered arousal is not target-specific editing.

The causal effect of a target dimension $k$ on an outcome $\phi$ is written
\begin{equation}
\tau_k=\mathbb{E}[\phi\mid do(u_k=1),c]-\mathbb{E}[\phi\mid do(u_k=0),c],
\label{eq:causal}
\end{equation}
where $u_k$ is a targeted representation intervention and $c$ fixes stimulus, task, state, nuisance variables, and recipient context. The intervention is promoted only if the target effect generalizes to held-out cues, exceeds sham and non-specific disruption, and leaves non-target endpoints within preregistered bounds.

\subsection{D5: Cross-session and cross-subject transfer}
Transfer is tested in increasing strength. First, a classifier or retrieval policy is transferred across sessions. Next, a task-defined association is transferred across models or subjects. Then, a digital latent representation is loaded into an independent recipient model and tested for neural-state and behavioral alignment. A future biological recipient study would be the strongest test, but is outside the present execution and requires separate approvals.

Representation alignment is learned only on training data. Baselines include subject shuffling, non-target representations, and semantic-similarity alignment. Evaluation separates content, temporal order, relations, neural reinstatement, behavior, and collateral effects. A recipient that reproduces a motor response but fails relational probes is classified as policy transfer, not memory transplantation. Two participants may share task structure without sharing one subjective episode; the protocol preserves this distinction.

\subsection{D6: Promotion decision}
The promotion rule is conjunctive. A level is promoted only when (i) the target content or relation was specified before evaluation, (ii) the effect generalizes to held-out cues or contexts, (iii) non-target and nuisance endpoints remain within bounds, (iv) the effect exceeds sham and matched non-specific disruption, and (v) an independent content scorer agrees with neural and behavioral evidence. Thus, classifier success cannot promote L0 to L3, neural reinstatement cannot automatically promote L2 to L3, and animal fear expression cannot promote any result to human subjective recollection.

\begin{table}[t]
\caption{Fidelity gates and failure interpretations}
\label{tab:gates}
\centering
\footnotesize
\setlength{\tabcolsep}{2pt}
\begin{tabular}{p{0.13\columnwidth}p{0.30\columnwidth}p{0.43\columnwidth}}
\toprule
Level & Required positive evidence & Interpretation if the next gate fails \\
\midrule
L0 & Held-out prediction of item or recall & Predictive access only. \\
L1 & Target-specific constrained behavior & Association or policy, not content storage. \\
L2 & Held-out neural reinstatement & Neural correlate without automatic content completeness. \\
L3 & Content, order, relations, bounded errors & Reconstruction within the tested phenotype. \\
L4 & Independent recipient passes all tests & Transfer claim only if identity and selectivity are justified. \\
\bottomrule
\end{tabular}
\end{table}

\section{Expected Outcome Branches and Conditional Conclusions}
The design is intended to produce informative branches rather than one headline number. If L0 succeeds but L1 fails, the system has predictive access without a reliable intervention. If L1 succeeds but L2 fails, a behavioral policy has transferred without neural reinstatement. If L2 succeeds but L3 fails, a neural correlate has been reproduced while content remains incomplete or confabulated. If L3 succeeds within a subject but L4 fails across subjects, digital storage and reconstruction may be feasible for a defined phenotype while transplantation is not established. If false details increase after manipulation, the fidelity gate fails even if recognition improves.

A strong positive result would be a selective, held-out, cross-context reconstruction of content and relations with calibrated uncertainty, low false-detail rate, and a causal perturbation effect that is specific to the target. Such a result would justify the claim that a task-defined memory representation is functionally transferable at the demonstrated level. It would not justify the claim that an entire autobiographical self has been uploaded. A partial positive result is scientifically valuable: for example, preserving temporal order and relations while losing affective nuance would identify a concrete boundary for the next design.

The most important negative result would also be informative. If performance collapses under novel cue combinations, cross-session transfer, or nuisance-matched controls, the apparent memory signal may be stimulus, strategy, or state decoding. If confidence and narrative fluency rise with false details, the system may be optimizing plausibility rather than memory. If target edits alter neighboring memories or agency, selective manipulation has failed even when the target score improves. These branches turn speculative language into a cumulative test program.

The near-term answer to the motivating question is therefore qualified but direct. Human memory can plausibly be stored digitally in progressively useful, task-defined representations, and current neuroscience already supports causal manipulation of memory-related processes in limited settings. Complete digital transplantation of human episodic recollection is not established. The decisive experiment is not a larger decoder; it is an independently validated, selective, content-preserving transfer test with explicit failure accounting.

\section{Risks, Limitations, and Human Review Requirements}
\subsection{Scientific and technical risks}
The representation may encode generic arousal, movement, task structure, semantic similarity, or subject-specific strategy. Cross-validation leakage can make a decoder appear to preserve content. Neural pattern similarity can reflect shared preprocessing or a common cue rather than reinstated memory. Generative models can insert details that were never observed. These risks are addressed by subject and session splits, nuisance-matched controls, novel relational probes, shuffled baselines, uncertainty reporting, and independent scoring.

The phenotype itself can be too narrow. A system may pass a paired-associate task and fail autobiographical context, or pass item identity while losing temporal order and source. The paper therefore does not compress all endpoints into one universal score. Its content score in (\ref{eq:content}) is accompanied by its components and is interpreted only for the declared task. Subjective recollection remains an endpoint requiring careful psychophysical and source-monitoring validation, not a quantity inferred from neural similarity alone.

\subsection{Ethics, governance, and autonomy}
Human neural data require consent, privacy protection, access governance, and controls over re-identification. Any human intervention requires institutional approval, safety monitoring, and a clear separation between research participation and clinical benefit. Memory alteration may affect identity, relationships, responsibility, trauma, and decision autonomy. Agency, confidence, affect, and neighboring memories must be measured rather than assumed. Participants must be able to understand and refuse the specific manipulation, and the protocol must not treat an improved task score as evidence that an intervention is benign.

Animal engram work requires animal-welfare review and a justification that the mechanistic question cannot be answered with a lower-burden design. It should not be used to imply human subjective experience. Future biological recipient studies require additional review for consent, recipient welfare, identity claims, and reversibility. The present artifact contains no executable stimulation recipe and no authorization to perform any neural manipulation.

\subsection{Human review and reproducibility}
Before execution, reviewers should inspect data-use permissions, consent, animal or human-subject approvals, power and sample assumptions, model leakage, content rubrics, false-memory risk, stimulation safety, and autonomy safeguards. The versioned object should retain data identifiers, preprocessing, model versions, fold assignments, random seeds, perturbation definitions, independent scoring rubrics, null results, unknowns, and all manual decisions. A failed language-model batch or an incomplete label must be preserved as a failure record and marked \texttt{needs\_human\_input}; it must not be silently repaired into an observed result.

\section{Conclusion}
The upload metaphor asks for a yes-or-no answer to a layered scientific problem. DEFL-CTA replaces that binary with a cumulative standard. Digital memory research begins with prediction and association, advances through neural reinstatement and content-preserving reconstruction, and reaches transfer only when an independent recipient passes held-out content, relational, neural, behavioral, and collateral tests. The ladder is demanding because memory is distributed, dynamic, and reconstructive, not because digital representation is inherently impossible.

The proposal makes room for real progress now. A versioned task-defined representation could support memory prostheses, experimental models of reconsolidation, and clinically relevant measurement of false detail without pretending to upload a person. At the same time, the program identifies a clear frontier: selective causal transfer of a declared memory phenotype across contexts and, eventually, across independent recipients. Until that frontier is met, the scientifically correct statement is not that human memory has been transplanted digitally, but that specific components of memory-related information may be decoded, modeled, and manipulated under controlled conditions. DEFL-CTA provides the audit trail needed to determine how far a future system has actually traveled.

\appendices
\section{Fidelity Ladder and Operational Definition Appendix}
Table~\ref{tab:appendix_ladder} gives the operational separation used by the Author Agent. The table is intentionally conjunctive: a higher level inherits the lower-level measurement burden and adds a new type of evidence. The design does not allow a model to claim L4 by reporting only a neural or behavioral endpoint.

\begin{table}[h]
\caption{Operational fidelity ladder}
\label{tab:appendix_ladder}
\centering
\footnotesize
\setlength{\tabcolsep}{2pt}
\begin{tabular}{p{0.10\columnwidth}p{0.34\columnwidth}p{0.43\columnwidth}}
\toprule
Level & Object being tested & Minimum evidence \\
\midrule
L0 & Access to a target label & Held-out prediction with leakage and nuisance controls. \\
L1 & Target-related action or association & Selective behavior above matched controls. \\
L2 & Target-related neural state & Reinstatement under held-out cues and contexts. \\
L3 & Declared memory phenotype & Content, order, source, and relation recovery with error decomposition. \\
L4 & Independent recipient expression & L3 content plus neural and behavioral alignment, selectivity, and independent validation. \\
\bottomrule
\end{tabular}
\end{table}

The ontology records seven separable fields: item or event identity; temporal order; source context; relational structure; neural reinstatement; behavioral expression; and subjective recollection. It also records omission rate, false-detail rate, confidence calibration, collateral-memory effect, agency, and autonomy. The fields may be correlated, but no field is used as a substitute for another. In particular, behavioral expression is not subjective recollection, and neural reinstatement is not content-complete reconstruction.

\section{Evidence Coverage and Review Appendix}
The Survey-to-Author binding uses project identifier \texttt{neur\_3}, Survey run identifier \texttt{neur\_3-survey-001}, and project-context fingerprint \texttt{neur3-digital-memory-context-v1}. The canonical artifacts are the Survey evidence plan, gap ledger, Idea result, ExperimentDesign JSON, and Author handoff stored with this report. The accepted gaps are G1--G6; G7 and G8 remain candidate gaps requiring validation. The source registry is E1--E10.

Claims about distributed and reconstructive memory are anchored to E1, E2, E8, and E9. Claims about rodent engram manipulation are bounded by E2--E5. Claims about human neural assistance are bounded by E6 and E7. Ethical requirements are anchored to E9 and E10. The document status is \texttt{PROPOSAL\_NO\_OBSERVED\_RESULTS}; observed results remain an empty set. No chapter claims a completed upload, digital edit, human memory transplant, or independently proven subjective recollection.

\begin{thebibliography}{10}
\bibitem{Squire2011} L. R. Squire and S. M. Z. Wixted, ``The cognitive neuroscience of human memory since H.M.,'' \textit{Neuron}, vol. 70, no. 3, pp. 435--445, 2011.
\bibitem{Josselyn2020} S. A. Josselyn and S. Tonegawa, ``Memory engrams: Recalling the past and imagining the future,'' \textit{Science}, vol. 367, no. 6473, 2020.
\bibitem{Liu2012} X. Liu \textit{et al.}, ``Optogenetic stimulation of a hippocampal engram activates fear memory recall,'' \textit{Nature}, vol. 484, pp. 381--385, 2012.
\bibitem{Ramirez2013} S. Ramirez \textit{et al.}, ``Creating a false memory in the hippocampus,'' \textit{Science}, vol. 341, no. 6144, pp. 387--391, 2013.
\bibitem{Chen2019} B. K. Chen \textit{et al.}, ``Artificially enhancing and suppressing hippocampus-mediated memories,'' \textit{Current Biology}, vol. 29, no. 11, pp. 1885--1894.e4, 2019, doi: 10.1016/j.cub.2019.04.065.
\bibitem{Hampson2018} R. E. Hampson \textit{et al.}, ``Developing a hippocampal neural prosthetic to facilitate human memory encoding and recall,'' \textit{Journal of Neural Engineering}, vol. 15, no. 3, 2018.
\bibitem{Ezzyat2018} M. Ezzyat \textit{et al.}, ``Closed-loop stimulation of temporal cortex rescues functional networks and improves memory,'' \textit{Nature Communications}, vol. 9, art. no. 3658, 2018.
\bibitem{Kandel2014} E. R. Kandel, Y. Dudai, and M. R. Mayford, ``The molecular and systems biology of memory,'' \textit{Cell}, vol. 157, no. 1, pp. 163--186, 2014.
\bibitem{Schacter1999} D. L. Schacter, ``The seven sins of memory: Insights from psychology and cognitive neuroscience,'' \textit{American Psychologist}, vol. 54, no. 3, pp. 182--203, 1999.
\bibitem{Ienca2017} M. Ienca and R. Andorno, ``Towards new human rights in the age of neuroscience and neurotechnology,'' \textit{Life Sciences, Society and Policy}, vol. 13, art. no. 5, 2017.
\end{thebibliography}

\end{document}
